\subsection{GROUPS WHOSE LIE ALGEBRA IS COMPACT}

THEOREM 1. (H. Weyl) Let G be a connected Lie group whose Lie algebra is compact semi-simple. Then G is compact and its centre is finite.

Since G is semi-simple, its centre D is discrete. Moreover, the quotient group G/D is isomorphic to \(\operatorname{Ad}(G)\) (Chap. III, §6, no. 4, Cor. 4), hence compact (Prop. 3). Finally, the group G/D is equal to its derived group (Chap. III, §9, no. 2, Cor. of Prop. 4). The theorem now follows from Integration, Chap. VII, §3, no. 2, Prop. 5.

PROPOSITION 4. Let G be a connected Lie group whose Lie algebra is compact. There exist a torus T, a simply-connected compact semi-simple Lie group S, a vector group V and a surjective morphism \(f: V \times T \times S \to G\) with finite kernel. If G is compact, the group V is reduced to the identity element.

Let C (resp. S) be a simply-connected Lie group whose Lie algebra is isomorphic to the centre (resp. the derived algebra) of L(G). Then C is a vector group, S is a compact group with finite centre (Th. 1) and G can be identified with the quotient of C × S by a discrete subgroup D, which is central (Integration, Chap. VII, §3, no. 2, Lemma 4). Since the image of the projection of D onto S is central, hence finite, D∩C is of finite index in D. Let C' be the vector subspace of C generated by D ∩ C, and V a complementary subspace. Then the group T = C'/(D ∩ C) is a torus, and G is isomorphic to the quotient of the product group V × T × S by a finite group.

If G is compact, so is \(V \times T \times S\) (General Topology, Chap. III, §4, no. 1, Cor. 2 of Prop. 2), hence so is V, which implies that \(V = \{e\}\) .

COROLLARY 1. Let G be a connected compact Lie group. Then \(\mathrm{C}(\mathrm{G})_{0}\) is a torus, \(\mathrm{D}(\mathrm{G})\) is a connected compact semi-simple Lie group and the morphism \((x,y)\mapsto xy\) from \(\mathrm{C}(\mathrm{G})_{0}\times\mathrm{D}(\mathrm{G})\) to G is a finite covering.

With the notation in Prop. 4, we have \(\mathrm{V} = \{e\}\) and the subgroups \(f(\mathrm{T})\) and \(f(\mathrm{S})\) of \(\mathrm{G}\) are compact, hence closed. Thus it suffices to show that \(f(\mathrm{T}) = \mathrm{C}(\mathrm{G})_0\) , \(f(\mathrm{S}) = \mathrm{D}(\mathrm{G})\) . Now, \(\mathrm{L}(\mathrm{G}) = \mathrm{L}(f(\mathrm{T})) \times \mathrm{L}(f(\mathrm{S}))\) ; since \(\mathrm{S}\) is semi-simple and \(\mathrm{T}\) is commutative, this implies that \(\mathrm{L}(f(\mathrm{T})) = \mathcal{C}(\mathrm{L}(\mathrm{G})) = \mathrm{L}(\mathrm{C}(\mathrm{G})_0)\) (Chap. III, §9, no. 3, Prop. 8) and \(\mathrm{L}(f(\mathrm{S})) = \mathcal{DL}(\mathrm{G}) = \mathrm{L}(\mathrm{D}(\mathrm{G}))\) (Chap. III, §9, no. 2, Cor. of Prop. 4), hence the stated assertion.

COROLLARY 2. The centre and the fundamental group of a connected compact semi-simple Lie group are finite. Its universal covering is compact.

With the notation in Prop. 4, the groups V and T are reduced to the identity element; thus S is a universal covering of G, and the fundamental group of G is isomorphic to \(\operatorname{Ker} f\) , hence finite. The centre D of G is discrete since G is semi-simple, so D is finite.

COROLLARY 3. The fundamental group of a connected compact Lie group G is a Z-module of finite type, of rank equal to the dimension of C(G).

Indeed, with the notations in Cor. 1, the fundamental group of \(\mathrm{C}(\mathrm{G})_{0}\) is isomorphic to \(Z^{n}\) , with \(n = \dim \mathrm{C}(\mathrm{G})_{0}\) , and the fundamental group of \(\mathrm{D}(\mathrm{G})\) is finite (Cor. 2).

COROLLARY 4. Let G be a connected compact Lie group. The following conditions are equivalent:

\begin{enumerate}
\def\labelenumi{(\roman{enumi})}
\item
  G is semi-simple;
\item
  C(G) is finite;
\item
  \(\pi_{1}(G)\) is finite.
\end{enumerate}

If \(\mathrm{G}\) is simply-connected, it is semi-simple.

This follows from Cor. 1 to 3.

COROLLARY 5. Let G be a connected compact Lie group. Then Int(G) is the identity component of the Lie group Aut(G) (Chap. III, §10, no. 2).

Let \(f \in \mathrm{Aut}(\mathrm{G})_0\) . Then \(f\) induces an automorphism \(f_1\) of \(\mathrm{C}(\mathrm{G})_0\) and an automorphism \(f_2\) of \(\mathrm{D}(\mathrm{G})\) , and we have \(f_1 \in \mathrm{Aut}(\mathrm{C}(\mathrm{G})_0)_0\) , \(f_2 \in \mathrm{Aut}(\mathrm{D}(\mathrm{G}))_0\) . Since \(\mathrm{Aut}(\mathrm{C}(\mathrm{G})_0)\) is discrete (General Topology, Chap. VII, §2, no. 4, Prop. 5), we have \(f_1 = \mathrm{Id}\) ; since \(\mathrm{D}(\mathrm{G})\) is semi-simple, by Chap. III, §10, no. 2, Cor. 2 of Th. 1 there exists an element \(g\) of \(\mathrm{D}(\mathrm{G})\) such that \(f_2(x) = g x g^{-1}\) for all \(x \in \mathrm{D}(\mathrm{G})\) . For all \(x \in \mathrm{C}(\mathrm{G})_0\) , we have \(g x g^{-1} = x = f_1(x)\) ; since \(\mathrm{G} = \mathrm{C}(\mathrm{G})_0.\mathrm{D}(\mathrm{G})\) , it follows that \(g x g^{-1} = f(x)\) for all \(x \in \mathrm{G}\) , so \(f = \operatorname{Int} g\) .

PROPOSITION 5. Let \(\mathbf{G}\) be a Lie group whose Lie algebra is compact.

\begin{enumerate}
\def\labelenumi{\alph{enumi})}
\item
  Assume that G is connected. Then G has a largest compact subgroup K; it is connected. There exists a closed central vector subgroup (no. 2) N of G such that G is the direct product \(N \times K\) .
\item
  Assume that the group of connected components of G is finite. Then:
\end{enumerate}

\begin{enumerate}
\def\labelenumi{(\roman{enumi})}
\item
  Every compact subgroup of \(\mathbf{G}\) is contained in a maximal compact subgroup.
\item
  If \(\mathrm{K}_1\) and \(\mathrm{K}_2\) are two maximal compact subgroups of \(\mathbf{G}\) , there exists \(g \in \mathbf{G}\) such that \(\mathrm{K}_2 = g\mathrm{K}_1g^{-1}\) .
\item
  Let \(\mathrm{K}\) be a maximal compact subgroup of \(\mathrm{G}\) . Then \(\mathrm{K} \cap \mathrm{G}_0\) is equal to \(\mathrm{K}_0\) ; it is the largest compact subgroup of \(\mathrm{G}_0\) .
\item
  There exists a closed central vector subgroup N of \(G_{0}\) , normal in G, such that, for any maximal compact subgroup K of G, \(G_{0}\) is the direct product of \(K_{0}\) by N, and G is the semi-direct product of K by N.
\end{enumerate}

\begin{enumerate}
\def\labelenumi{\alph{enumi})}
\item
  We retain the notations of Prop. 4. The projection of \(\mathrm{Ker}f\) onto V is a finite subgroup of the vector group V, hence is reduced to the identity element. It follows that \(\mathrm{Ker}f\) is contained in \(\mathrm{T} \times \mathrm{S}\) , hence that G is the direct product of the vector group \(\mathrm{N} = f(\mathrm{V})\) and the compact group \(\mathrm{K} = f(\mathrm{T} \times \mathrm{S})\) . Every compact subgroup of G has a projection onto N that is reduced to the identity element, hence is contained in K. This proves \(a\) ).
\item
  Assume now that \(G/G_{0}\) is finite. By a), \(G_{0}\) is the direct product of its largest compact subgroup M and a vector subgroup P; the subgroup M of G is clearly normal. Let n be a vector subspace complement of \(L(M)\) in \(L(G)\) , stable under the adjoint representation of G (no. 1 and no. 3, Prop. 3); this is an ideal of \(L(G)\) and we have \(L(G) = L(M) \times n\) . Let N be the integral subgroup of G with Lie algebra n; by Chap. III, §6, no. 6, Prop. 14, it is normal in G. The projection of \(L(G)\) onto \(L(P)\) with kernel \(L(M)\) induces an isomorphism from n to \(L(P)\) ; it follows that the projection of \(G_{0}\) onto P induces an étale morphism from N to P; since P is simply-connected, this is an isomorphism, and N is a vector group. The morphism \((x, y) \mapsto xy\) from \(M \times N\) to \(G_{0}\) is an injective étale morphism (since \(M \cap N\) is reduced to the identity element), hence an isomorphism. It follows that N is a closed subgroup of G and that the quotient G/N is compact, since \(G_{0}/N\) is compact and \(G/G_{0}\) is finite (General Topology, Chap. III, §4, no. 1, Cor. 2 of Prop. 2).
\end{enumerate}

By Integration, Chap. VII, §3, no. 2, Prop. 3, every compact subgroup of G is contained in a maximal compact subgroup, these are conjugate, and for any maximal compact subgroup K of G, G is the semi-direct product of K by N. Since \(G_{0}\) contains N, it is the semi-direct product of N by \(G_{0} \cap K\) ; it follows that \(G_{0} \cap K\) is connected, hence equal to \(K_{0}\) , since \(\mathrm{K}/(\mathrm{G}_{0} \cap \mathrm{K})\) is isomorphic to \(G/G_{0}\) , hence finite; finally, \(K_{0}\) is clearly the largest compact subgroup of \(G_{0}\) by a).

COROLLARY. If N satisfies the conditions of b) (iv), and if \(K_{1}\) and \(K_{2}\) are two maximal compact subgroups of G, there exists \(n \in N\) such that \(nK_{1}n^{-1} = K_{2}\) .

Indeed, by (ii) there exists an element \(g \in G\) such that \(gK_{1}g^{-1} = K_{2}\) ; by (iv), there exists \(n \in N\) and \(k \in K_{1}\) such that g = nk. The element n then has the required properties.

